\section{Introduction}

By 2022, more than 2,000 \acp{bss} were available worldwide, with approximately 9 million bikes in circulation~\cite{MeddinWorldMap2025}. These systems enhance the efficiency of \ac{pt} by improving access to stations~\cite{Noland2019, Kapuku2021}, which in turn reduces car usage and alleviates traffic congestion~\cite{Ma2020, Fan2020}. Moreover, the increased use of bikes and the resulting shift in transport modes have been shown to positively impact public health by promoting physical activity~\cite{Otero2018}, benefit the environment by reducing pollutant emissions~\cite{Clockston2021}, and lead to economic savings~\cite{Ricci2015}.

The recent increase in data availability related to \acp{bss} has led to greater interest in understanding the mobility patterns within these systems~\cite{Kou2019, Li2020}. Studies have shown that stations with strong flows tend to be located nearby~\cite{Chen2017}, and they often display a certain level of mesoscopic organization within communities~\cite{Zaltz2013}. Additionally, the development of models to understand the usage of micro-mobility~\cite{Tran2015} has improved understanding of factors such as topography~\cite{Kim2020}, day of the week~\cite{Tran2015}, weather conditions~\cite{ElAssi2017, Kim2018, Ashqar2019}, and the built environment~\cite{Tran2015}. Furthermore, efforts have been made to develop predictive models that enhance the management of \acp{bss} infrastructures efficiently and guide system expansion through \ac{ml} and probabilistic approaches~\cite{Singhvi2015, Li2015, Chen2016, Albuquerque2021, Yang2020graph, Almannaa2020}. In this line, one of the major challenges in optimizing \ac{bss} performance is ensuring bike availability, as stations often experience significant flow imbalances~\cite{Yang2016, Singla2015, FaghihImani2017Empirical}. To address this, models focused on optimizing rebalancing operations have been developed~\cite{Raviv2013, Dell2014, Dell2018, Chiariotti2018}.

In recent years, \acp{bss} have undergone progressive electrification~\cite{MeddinWorldMap2025, Galatoulas2020}, with a sharp increase in the number of \acp{e-b} compared to their mechanical counterparts~\cite{Galatoulas2020, Bielinski2020}, as \acp{e-b} allow for longer trips and attract a broader user base~\cite{Simsekoglu2019, Weinert2007}. \acp{e-b} have been shown to be less sensitive to factors such as longer distances, adverse weather conditions, and casualties~\cite{Wamburu2021, Campbell2016, Siman2018}. Despite these advantages, \acp{e-b} present additional challenges to maintaining bike availability, primarily due to low battery levels, battery malfunctions, and complications arising from high temperatures~\cite{Ji2014, Florez2018, Zhang2019, Usama2019}. Moreover, the scarcity of battery data at the system level complicates the development of predictive models. While several studies have implemented frameworks to predict the battery consumption of individual bikes during trips~\cite{Burani2022, Steyn2014}, most approaches at the \ac{bss} level focus on the station-level granularity~\cite{Zhou2023}.

This publication addresses the challenge that the state of charge poses to \ac{e-b} availability by analyzing mobility and battery level data from Barcelona's \ac{bss}, \textit{Bicing}, and proposing a model that predicts future bike locations with their corresponding battery levels. The structure of this Chapter is as follows: first, the spatial and temporal variability in average battery levels is analyzed to identify regions more vulnerable to battery depletion. Next, a Markov chain approach is implemented to predict the future locations of bikes. Subsequently, the parameters in the equations describing battery power consumption are optimized. Finally, by combining the models presented in the two previous Sections, a framework is provided to support the management of \ac{e-b} fleets by projecting future scenarios.